\documentclass[11pt]{article}
\usepackage[margin=1in]{geometry}
\usepackage{amsmath,amssymb}
\usepackage[hidelinks]{hyperref}

\newcommand{\ellp}{\ell^{p(\cdot)}}
\newcommand{\cnull}{c_0}
\newcommand{\norm}[1]{\left\|#1\right\|}
\newcommand{\N}{\mathbb N}

\begin{document}

\section*{Human Review: Variable-Exponent \texorpdfstring{$c_0$}{c0} Tail Counterexample}

\noindent Original automatic proof: \texttt{solution\_packet.pdf}

\noindent Checked date: 15 June 2026

\noindent Status: Manually checked.

\noindent Verdict: The proof appears correct.

\section*{Human Comments}

The proof appears to give a correct negative answer to Talponen's question asking whether
\[
  \ell^{p(\cdot)}\cap c_0
\]
always coincides with the closed linear span
\[
  [(e_n)]\subset \ell^{p(\cdot)}.
\]
Here \([(e_n)]\) means the closed linear span of the canonical unit vectors with respect to the \(\ell^{p(\cdot)}\)-norm, not the \(c_0\)- or \(\ell_\infty\)-norm.

The construction is quite nice. The vector is built in blocks. In block \(j\), the exponent is large,
\[
  r_j=j+1,
\]
and the block contains
\[
  m_j=j^{r_j}
\]
coordinates, each equal to \(c/j\), where \(0<c<1\). Since \(c/j\to0\), the resulting vector belongs to \(c_0\).

The key idea is that each block has two different behaviours depending on how it is viewed. If the block is viewed by itself, starting from zero, its norm is exactly
\[
  \left(m_j\left(\frac cj\right)^{r_j}\right)^{1/r_j}=c.
\]
Thus tails beginning just before block \(j\) always contain a finite piece of norm \(c\). Consequently the tails cannot converge to zero in the \(\ell^{p(\cdot)}\)-norm, and this prevents the vector from lying in \([(e_n)]\).

On the other hand, when the same block is appended after the previous nonzero part of the vector, the recursive norm gives
\[
  A_j=\left(A_{j-1}^{r_j}+c^{r_j}\right)^{1/r_j}.
\]
Since \(r_j\to\infty\) and \(0<c<1\), the extra contribution \(c^{r_j}\) becomes very small. The estimate
\[
  A_j\leq A_{j-1}(1+c^{r_j})^{1/r_j}
\]
leads to the convergent product
\[
  \prod_j (1+c^{r_j})^{1/r_j},
\]
because
\[
  \sum_j \frac{c^{r_j}}{r_j}<\infty.
\]
This shows that the norms of the finite initial segments are uniformly bounded, hence the vector belongs to \(\ell^{p(\cdot)}\).

The auxiliary tail characterization also appears correct:
\[
  y\in[(e_n)]
  \quad\Longleftrightarrow\quad
  \norm{Q_ny}_{\ell^{p(\cdot)}}\to0.
\]
The proof uses only approximation by finitely supported vectors and contractivity of coordinate projections. The contractivity follows from coordinatewise monotonicity of the recursive norm, since coordinate projections only decrease the absolute values of coordinates.

\section*{Review Outcome}

Archive this packet as an apparently correct counterexample. The proof's main mechanism is sound: each block remains large when seen as a fresh tail, but contributes less and less when appended after the previously accumulated part of the vector. This gives
\[
  x\in\ell^{p(\cdot)}\cap c_0
  \qquad\text{but}\qquad
  x\notin[(e_n)].
\]
I would recommend checking with the author whether the question has since been answered elsewhere, but mathematically the proposed construction appears correct.

\end{document}
